\label{app:implementation}

Concept-to-code map and engineering notes for the search engine of
Section~\ref{sec:search}.

\paragraph{Concept-to-class map.} The main concepts of
Section~\ref{sec:search} correspond to the following classes and
flags in the reference implementation:

\begin{center}
\begin{tabular}{ll}
\toprule
concept & class / flag \\
\midrule
frontier driver (impact-DAG priority queue) & \code{FrontierClosure} \\
recombination materialiser & \code{Recombination.constructWithAllocation} \\
composed sub-product lookup (overlay + disk) & \code{AlgorithmLookup}, \code{FieldAwareLookup} \\
allocation branch-and-bound & \code{AllocationOptimizer} \\
search budget & \code{SearchBudget}: \code{EXACT}, \code{upTo}, \code{maxNodes} \\
allocation filters & \code{maxImbalance}, \code{maxCombinations} \\
axis-flip orbit mask evaluation & \code{AnalyticalMaskSearch} \\
seed ranking by \(\omega\) & \code{docs/bases-by-omega.md} \\
\bottomrule
\end{tabular}
\end{center}

\paragraph{Pool configuration.} The default pool
(\code{PoolConfig.simple}) carries a small number of high-value outer
schemes --- Strassen, Winograd, Smirnov 3-row templates, and a few
imported Pan / AT / FMM bases. An \code{extendedPool} mode adds every
\code{Leaf}-tagged catalog entry as an outer template; that pool is
large enough that the \code{maxCombinations}-capped allocation
enumerator becomes load-bearing.

\paragraph{Budget flags.} \code{AllocationOptimizer} runs under a
\code{SearchBudget}: \code{EXACT} (no cap), \code{upTo} a known upper
bound (e.g.\ the Kronecker rank), or a \code{maxNodes} anytime cap.
The result carries an \code{exhaustive} flag --- propagated to the
optimality tier of Section~\ref{sec:intro} --- and reports the
\code{nodes} visited against the \code{fullSpace} of allocations.

\paragraph{Seed ranking.} \code{docs/bases-by-omega.md} is an
automatically-regenerated table ranking every pool-eligible base by
its raw \(\omega\), its \(\omega\) after imbalance-\(\le 3\)
restriction, and its \(\omega\) after axis-flip-canonical and
S\(_3\)-canonical orbit reductions, evaluated at the reference target
\(\nmpshape{17}{17}{17}\).

\paragraph{Caching.} Caching exists at two layers. The first is
per-scheme \emph{symbolic templates}: for each outer scheme we extract
once a tuple of six block-support bitmasks per product (\(U\)-row,
\(U\)-col, \(V\)-row, \(V\)-col, \(W\)-row, \(W\)-col), and per-call
sub-shape extraction reduces from \(O(r \cdot \text{dim}^2)\) to
\(O(r \cdot \text{dim})\). The second is a per-(scheme, allocation,
peel) cache of the resulting sub-shape multiset; this is a
hashing-overhead trade and is currently capped at 200k entries per
scheme. A symbolic-template-at-allocation-pattern layer would
generalise the second cache; the current implementation captures the
common cases without paying the templating cost.

\paragraph{Progress.} A derivation closure over the entire catalog up to
dim 32 takes hours of wall-clock; long runs emit periodic progress
lines and ETAs at three nested levels (per-round, per-shape, and
intra-search).
